% !TeX spellcheck = en_GB
% Template for Federated Learning Course Project Report
% Based on the ICASSP 2026 conference template (spconf.sty)

\documentclass{article}
\usepackage{spconf,amsmath,graphicx,hyperref}
\usepackage{bm,amsfonts,amssymb,amsbsy}
\usepackage[dvipsnames]{xcolor}
\usepackage{cleveref}
\usepackage{algorithm, algorithmic}
\usepackage{subcaption, adjustbox}
\usepackage{listings}
\usepackage{qrcode}

\input{ml_macros.tex}

% Minimalistic hint environment
\newenvironment{hint}{%
	\par\medskip\noindent
	\begin{minipage}[t]{\dimexpr\columnwidth-8mm}%
	\raggedright\itshape
	\textbf{Hint (remove before submission):}\space
}{%
	\end{minipage}\par\medskip
}

% Title
\title{Template for a Federated Learning Project Report}
\name{No Name (to make peer grading double-blind!)}
\address{}

\begin{document}

\maketitle

\begin{hint}
	All hints in this template are instructions for drafting and must be
	removed before final submission. Not every hint corresponds to a
	separate grading criterion; however, ignoring hints typically leads to unclear
	problem formulations, weak experiments, or low review scores.
\end{hint}

\begin{abstract}
	\begin{hint}
The abstract provides a concise summary of the project, including the FL application,
empirical graph modelling, variation minimization approach, and the FL algorithms used.
Do not include detailed numerical results. Instead, focus on problem, method, and evaluation setup.
	\end{hint}
\end{abstract}

\begin{keywords}
Federated learning, networks, personalized machine learning, trustworthy artificial intelligence
\end{keywords}

\begin{hint}
	Your report must adhere to the following formatting and length requirements:
\begin{itemize}
	\item Do not include any information (GitHub usernames, repository links, or dataset paths)
	           in the report that could be used to identify you.
	\item Use the \LaTeX\ settings as defined in this file without changing margins or font sizes.
	\item You must not modify the number or titles of sections in this template.
	\item The report must not exceed 5 pages in total. The 5\textsuperscript{th} page must
	contain only the references.
\end{itemize}
\end{hint}

\begin{hint}
You are encouraged to use terminology and notation consistent with the
Aalto Dictionary of Machine Learning \cite{AaltoDictML}.
The file \texttt{ml\_macros.tex} (included in this template) provides
macros such as \texttt{\textbackslash featurevec} for the feature vector,
\texttt{\textbackslash localparams\{i\}} for local model parameters at node $i$, etc.
Using these macros is not mandatory, but consistent notation
improves readability and peer review scores.
The dictionary repository also contains 300+ reusable TikZ figures
(FL networks, graphs, loss landscapes, etc.) that you can copy into your report.
\qrcode[height=1cm]{https://github.com/AaltoDictionaryofML/AaltoDictionaryofML.github.io}
\end{hint}


\section{Introduction}

\begin{hint}
	Describe the application scenario, motivation, and learning objective of your
	federated learning (FL) system.
	Your introduction should include:
	\begin{itemize}
		\item a clear \textbf{motivation} for why FL is suitable for your application,
		\item sufficient \textbf{background} on the problem domain and related work,
		\item a concise \textbf{outline} of the report structure (one sentence per section).
	\end{itemize}
\end{hint}


\section{Problem Formulation}
\label{sec:pf}
\begin{hint}
	All projects use the same raw data source and node definition:
	\begin{itemize}
		\item \textbf{Raw data:} Finnish Meteorological Institute (FMI) Open Data
		(\url{https://en.ilmatieteenlaitos.fi/open-data}).
		\item \textbf{Nodes (fixed):} each node in the FL network is one FMI weather station.
	\end{itemize}
	A Python script for fetching FMI data is available here:
	\qrcode[height=1cm]{https://github.com/alexjungaalto/FederatedLearning/blob/main/Edition2026/assets/GetFMIData.py}

	\medskip
	Everything else is \textbf{your design choice}. You must clearly
	describe and justify:
	\begin{itemize}
		\item which stations you selected and why,
		\item how edges between stations are defined and how edge weights are set
		(e.g., geographic proximity, data similarity);
		see \cite[Ch.~7]{Jung2025} for data-driven methods,
		\item how local data points are derived from raw FMI measurements
		(choice of features and labels),
		\item what model is trained locally at each station,
		\item how the local loss functions are defined.
	\end{itemize}
\end{hint}


\section{Methodology}
\label{sec:meth}
\begin{hint}
The project requires you to apply GTVMin-based methods to
the FL application modelled in Section~\ref{sec:pf}. In this section
you need to clearly state and explain:
\begin{itemize}
	\item Your choice of variation measure, e.g., $\gtvpenalty(\localparams{\nodeidx}-\localparams{\nodeidx'})$
	for parametric models.
	\item Your choice of FL algorithm (i.e., optimization method for solving
	GTVMin) and its synchronous message passing implementation.
\end{itemize}
You are not expected to propose novel FL algorithms; correct implementation and
analysis of existing GTVMin-based methods is sufficient.
\end{hint}


\section{Numerical Experiments}
\label{sec:experiments}

\begin{hint}
	Present quantitative results (training, validation, and test losses),
	comparisons, and figures.
	At a minimum, you must evaluate and compare \textbf{two distinct FL systems}.
	These systems must differ in a \emph{structural design choice}, such as:
	\begin{itemize}
		\item different graph constructions or edge weights,
		\item different local model classes or model capacities,
		\item different local loss functions.
	\end{itemize}
	Varying only hyper-parameters (e.g., regularization strength or step size)
	does \textbf{not} count as a distinct FL system.
	Clearly explain evaluation metrics (e.g., which loss function has been used
	for training, validation, and final evaluation on the test set) and experimental
	settings (e.g., how data is split into train/validation/test sets and how hyper-parameters
	are selected). Briefly explain how the nature of the data points was taken into
	account when choosing the splitting strategy (e.g., chronological splitting
	for time-series data).
\end{hint}

\begin{hint}
	\textbf{Figures and tables.}
	Figures should be designed to communicate quantitative results clearly and efficiently.
	Avoid decorative elements that do not convey information (e.g., excessive colours, 3D effects,
	or chart junk). Axes, legends, and captions must be self-contained and precise.
	Whenever possible, embed numerical values directly into figures or tables rather
	than relying on verbose textual explanations. For general principles of effective
	quantitative graphics, see \cite{TufteVDQI}.
\end{hint}

\begin{hint}
\textbf{No Python source code is required for submission.}
Instead, the report itself must provide sufficient detail so that the
FL algorithm can be \textbf{reimplemented using a large language model (LLM)}
such as Mistral Le Chat, Aleph Alpha, or DeepSeek via Ollama.
In particular, a reader should be able to paste the relevant sections
into an LLM and obtain working code. Ensure you specify:
\begin{itemize}
	\item \textbf{Data:} exact FMI station IDs (or names), time range, measured
	variables, and how features/labels are derived from raw measurements.
	\item \textbf{FL network:} which stations are connected and how edge weights
	are computed (with enough detail to recompute them).
	\item \textbf{Algorithm:} the update equation, learning rate, number of
	iterations, regularization parameter, and initialization.
\end{itemize}
\end{hint}



\section{Conclusion}
\label{sec:conclusion}
\begin{hint}
Interpret the numerical results from Sec.~\ref{sec:experiments}.
In particular,
\begin{itemize}
	\item Discuss whether the obtained results solve the problem satisfactorily.
	\item Identify limitations of the studied FL methods and suggest potential improvements (e.g.,
	collecting more data points or using larger local models at some nodes).
\end{itemize}
You can find some guidance on how to diagnose ML models in \cite[Sec.~6.6]{Jung2022}.
\end{hint}



\begin{thebibliography}{9}
	\bibitem{Jung2025} A.~Jung, \textit{Federated Learning: From Theory to Practice}, Springer, 2026. Arxiv Preprint: \url{https://arxiv.org/abs/2505.19183}.
	\bibitem{Jung2022} A.~Jung, \textit{Machine Learning: The Basics}, Springer, 2022. Available via Aalto Library. Preprint: \url{https://github.com/alexjungaalto/MachineLearningTheBasics/}.
	\bibitem{AaltoDictML} A.~Jung (Editor), \textit{Aalto Dictionary of Machine Learning}, \url{https://aaltodictionaryofml.github.io/}, 2026.
	\bibitem{TufteVDQI}
	E.~R.~Tufte,
	\emph{The Visual Display of Quantitative Information},
	2nd edition, Graphics Press, 2001.
\end{thebibliography}

\end{document}
